\subsection{The proper-tube theorem}

Definition~\ref{def:zero-object-hierarchy} reserves the word tube for a proper
submersion.  The resulting structure is strong and concrete.

\begin{theorem}[Proper real-zero tube]
\label{thm:proper-real-zero-tube}
Let $q:X\to B$ be an ambient regular AEG family over a connected one-dimensional
base, and suppose its assignment incidence $\calZ=A^{-1}(0)$ is nonempty and
$\pi=q|_{\calZ}$ is proper and surjective.
\begin{enumerate}
\item The map $\pi:\calZ\to B$ is a smooth fiber bundle whose fibers are smooth
one-manifolds.
\item If the surface fibers have no boundary, every zero fiber is a finite disjoint
union of circles.
\item If $B$ is an interval, the bundle is globally trivial.  For a product
ambient family $M\times B$, its compact zero curves over each compact subinterval
form a smooth isotopy and hence extend to an ambient isotopy of $M$.  If $B$ itself
is a compact interval, this applies over all of $B$.
\item If $B=S^1$, the surface fibers are boundaryless, and the smooth vertical
orientation required in Definition~\ref{def:ambient-aeg-family} is fixed, every
connected component of $\calZ$ is a torus.
\end{enumerate}
\end{theorem}

\begin{proof}
The AEG equation makes the vertical derivative $\dd A_b$ nonzero at every point,
so Theorem~\ref{thm:parameterized-zero-section} proves the bundle assertion.
Properness makes every fiber compact.  A compact boundaryless one-manifold is a
finite union of circles.

A smooth bundle over an interval is trivial.  In a product ambient family, a
trivialization expresses the inclusions of the compact fiber as a smooth isotopy of
embeddings over each compact subinterval; the isotopy-extension theorem extends it
to the ambient surface
\cite{Hirsch1976DifferentialTopology}.

Over $S^1$, the total space is the mapping torus of a diffeomorphism of a finite
union of circles.  Each orbit of the induced component permutation produces one
connected component.  After traversing the orbit, the return map is a
diffeomorphism of $S^1$, so that component is either a torus or a Klein bottle.
The smooth orientation of $T^VX$ and the orientation of the base orient $X$; the
regular level set $A^{-1}(0)$ is cooriented by $\dd A$, hence orientable.  The
Klein-bottle case is excluded, and every component is a torus.  This is where the
global vertical-orientation hypothesis is essential: a mapping torus with
orientation-reversing fiber monodromy need not be orientable.
\end{proof}

\begin{remark}[Why vertical orientability is necessary]
If the global vertical orientation is dropped, the torus conclusion is false even
for a product-type formula.  Let $M=\R_r\times S^1_\phi$ with
$g=\dd r^2+\dd\phi^2$, $a=\mu r$, and $\lambda=0$, and form the mapping torus of
the AEG isometry $(r,\phi)\mapsto(r,-\phi)$.  Its zero incidence is the mapping
torus of an orientation-reversing circle map, hence a Klein bottle.  It is compact
and proper over the parameter circle, but the resulting surface bundle has no
smooth orientation of $T^VX$.
\end{remark}

\begin{remark}[What this theorem does not encode]
The theorem classifies the abstract total surface under its hypotheses.  The
embedding $\calZ\hookrightarrow X$ may still carry nontrivial ambient monodromy,
and no preferred curve inside a torus component has been selected.  In particular,
a union of zero tori is not yet a knot or link.
\end{remark}

Example~\ref{ex:cylindrical-propagation} gives an immediate nonempty instance.  If
$B$ is compact and $c:B\to\R$ is smooth, set $a_b=a_{c(b)}$.  Then
\[
 \calZ=\{(r,\phi,b):r=c(b)\}\cong S^1\times B.
\]
It is compact and therefore proper over $B$.  The example is deliberately trivial:
it verifies every hypothesis while isolating the content supplied by properness.

\subsection{A compact helical AEG tube}

The next family is still explicit but has nontrivial component transport.  Let
\[
 M_L=[-L,L]_x\times S^1_y,
 \qquad B=S^1_\theta,
 \qquad L>0,
\]
and choose integers $n\ge1$ and $m\in\Z$.  Define
\begin{equation}
 A_{m,n}(x,y,\theta)=e^{nx}\sin(ny-m\theta)
 \label{eq:helical-assignment}
\end{equation}
and the vertical metric
\begin{equation}
 g_{m,n,\theta}=
 \frac{n^2e^{2nx}}
 {\mu^2+\lambda^2e^{2nx}\sin^2(ny-m\theta)}
 (\dd x^2+\dd y^2).
 \label{eq:helical-metric}
\end{equation}

\begin{theorem}[Helical zero-tube theorem]
\label{thm:helical-zero-tube}
The family \eqref{eq:helical-assignment}--\eqref{eq:helical-metric} is an ambient
regular AEG family and its zero incidence is a proper neat tube over $S^1$.
For every $\theta$, the zero fiber consists of $2n$ intervals
\begin{equation}
 y=\frac{m\theta+k\pi}{n},
 \qquad k\in\Z/(2n).
 \label{eq:helical-zero-fiber}
\end{equation}
Parameter monodromy acts on the component labels by
\begin{equation}
 k\longmapsto k+2m\pmod{2n}.
 \label{eq:helical-permutation}
\end{equation}
If $d=\gcd(n,m)$, including $d=n$ when $m=0$, the total incidence is a disjoint
union of $2d$ annuli, each containing $n/d$ zero intervals in its monodromy orbit.
\end{theorem}

\begin{proof}
The spatial derivative is
\[
 \dd^VA_{m,n}
 =ne^{nx}\sin(ny-m\theta)\,\dd x
  +ne^{nx}\cos(ny-m\theta)\,\dd y,
\]
whose Euclidean norm squared is $n^2e^{2nx}$.  Thus
\eqref{eq:helical-metric} is precisely the conformal realization metric, proving
the AEG equation.  At a zero, the $y$-derivative is nonzero, so the incidence is
smooth and meets $x=\pm L$ neatly.  It is closed in the compact manifold
$M_L\times S^1$, hence compact and proper.

Equation~\eqref{eq:helical-zero-fiber} follows from $ny-m\theta\in\pi\Z$.
Increasing $\theta$ by $2\pi$ changes the label by $2m$, proving
\eqref{eq:helical-permutation}.  Translation by $2m$ on the cyclic group of order
$2n$ has $\gcd(2n,2m)=2d$ orbits, each of length $n/d$.  The transport fixes the
$x$-coordinate, so the return map after a complete label orbit is the identity on
the interval.  Its mapping torus is therefore an annulus, not a M\"obius band.
\end{proof}

\begin{figure}[ht]
\centering
\begin{tikzpicture}[x=0.95cm,y=0.95cm,>=Latex]
  \draw[->] (0,0)--(7.1,0) node[right] {$\theta$};
  \draw[->] (0,-0.2)--(0,4.7) node[above] {$y$};
  \draw[dashed] (6,0)--(6,4.4);
  \node[below] at (0,0) {$0$};
  \node[below] at (6,0) {$2\pi$};
  \foreach \j in {-1,0,1,2,3} {
    \draw[blue!70!black,thick] (-0.2,{0.7+1.0*\j}) -- (6.2,{2.1+1.0*\j});
  }
  \draw[<->,red!70!black,thick] (6.35,0.7)--(6.35,2.7)
    node[midway,right,align=left,font=\small]{label shift\\$k\mapsto k+2m$};
  \node[align=center,font=\small] at (3.25,4.65)
    {unwrapped zero sheets $ny-m\theta=k\pi$};
\end{tikzpicture}
\caption{The helical tube on the unwrapped $(\theta,y)$-plane.  Vertical boundary
identification at $\theta=0,2\pi$ produces the permutation
\eqref{eq:helical-permutation}; the drawing is schematic, while the slope is fixed
exactly by $\dd y/\dd\theta=m/n$.}
\label{fig:helical-unwrapped}
\end{figure}

\subsection{Boundary traces and logarithmic sheets}

At either boundary component $x=\pm L$, the trace lies in the parameter-position
torus $S^1_\theta\times S^1_y$ and satisfies
\[
 ny-m\theta\in\pi\Z.
\]
It has $2d$ connected components.  Each is a primitive curve with homology class
\begin{equation}
 \frac nd[S^1_\theta]+\frac md[S^1_y].
 \label{eq:helical-boundary-class}
\end{equation}
If $T(p,q)$ denotes the torus link whose total homology data in these ordered
coordinates are $(p,q)$, the full trace is $T(2n,2m)$, up to simultaneous
orientation reversal.  This statement depends on the declared torus coordinates;
it is not an intrinsic knot type before an embedding of that torus in a
three-manifold is fixed.

Unwrap the $y$-circle to $y\in\R$.  Formula~\eqref{eq:helical-zero-fiber} now has
labels $k\in\Z$.  The parameter loop sends $k\mapsto k+2m$, while the deck
transformation $y\mapsto y+2\pi$ sends $k\mapsto k+2n$.  The finite downstairs
system is exactly the quotient of this countable label set by the deck action.
This simultaneously realizes a proper tube and a nontrivial logarithmic shift.

\begin{corollary}[Finite quotient of the helical lift]
\label{cor:helical-finite-quotient}
For the helical family, the countable lifted sheets contain no additional
component permutation beyond the affine action generated by the two translations
$2m$ and $2n$ on $\Z$.  Downstairs component monodromy is their reduction modulo
$2n$.
\end{corollary}

\subsection{Embedded monodromy versus braid monodromy}

The permutation \eqref{eq:helical-permutation} is honest monodromy, but it is not
an element of $B_{2n}$.  Its fiber consists of intervals in a surface, not points in
a disc.  Intersecting the zero tube with $x=c$ produces a one-dimensional trace in
$S^1_y\times S^1_\theta$, but the choice of $c$ is additional probe data.  Section
\ref{sec:threading} formalizes this distinction.

Even for compact circle fibers in Theorem~\ref{thm:proper-real-zero-tube}, the
abstract monodromy of the one-manifold fiber records only a permutation and
orientation-preserving return maps of circles.  The embedding in the ambient
surface can carry a richer mapping class, but that is a mapping-class problem for
curves, not the ordinary braid group of moving points.
